% !Mode:: "TeX:UTF-8"
\chapter{\label{chapter-yufa}\LaTeX{}基础语法}
\section{字体设置}
本模板使用了CTEX封装的\texttt{ctexbook}类，作为毕设论文，一般不需要过多使用字体。因此，全文默认字体为~{\bf 宋体}~和{\bf Times New Roman}字体。此外还可以使用\texttt{\textbackslash heiti}命令来使用~{\heiti 黑体}~字体。\par
对于英文，则有\textbackslash textbf\{•\}\quad\textbackslash textit\{•\}\quad\textbackslash textsf\{•\}\quad\textbackslash texttt\{•\}\quad\textbackslash textsl\{•\}命令来进行相应的控制。\par
{\textbackslash textbf\{•\}\qquad\textbf{0123456789abcdefgABCDEFG}}\par
{\textbackslash textit\{•\}\qquad\textit{0123456789abcdefgABCDEFG}}\par
{\textbackslash textsf\{•\}\qquad\textsf{0123456789abcdefgABCDEFG}}\par
{\textbackslash texttt\{•\}\qquad\texttt{0123456789abcdefgABCDEFG}}\par
{\textbackslash textsl\{•\}\qquad\textsl{0123456789abcdefgABCDEFG}}\par
在论文中，有一些地方需要进行强调说明时，可以使用~{\bf 加粗}~来进行强调，比如\{\textbackslash bf加粗\}，即可以得到~“{\bf 加粗}”~的效果；\par
还可以有~\underline{下划线}，使用命令为\textbackslash underline\{需要加下划线的内容\}即可；\par
也还可以通过对文字使用\textcolor{blue}{颜}\textcolor{red}{色}来强调说明。使用\textcolor[rgb]{0.5,0.95,0.96}{有颜色的文字}可以通过命令\texttt{\textbackslash textcolor}\{red\}\{红色的字\}，或者\texttt{\textbackslash textcolor[rgb]}\{0.5,0.98,0.7\}\{使用RGB混合颜色\}来\textcolor[rgb]{0.5,0.98,0.7}{使用RGB混合颜色}进行强调。
\section{字号设置}
论文要求正文字号为小四号，除非有特别的要求，全文不需要对字号进行另外的设置。需要进行字号设置时，可以用\texttt{\textbackslash zihao\{-4\}}，以及\texttt{\textbackslash zihao\{4\}}来表示小四号字体和四号字体，同理可以从小五号\texttt{\textbackslash zihao\{-5\}}字体到初号\texttt{\textbackslash zihao\{0\}}字体。\par
\section{间距间隔控制}
对于间距的控制，\LaTeX{}中的间距比Office要强大得多了！Office除了空格回车外，差不多就没有了。而对于\LaTeX{}来说，可使用的间隔控制命令有如表\ref{table-space}所示若干，虽然很多都用不到。\par
\begin{table}[ht]
\begin{center}
\caption{\LaTeX{}中的一些常用间隔间距控制命令\label{table-space}}
\begin{tabular}{l|l}
\hline
命令&效果\slash 说明\\\hline\hline
间隔\texttt{\textbackslash quad}间隔 & 间隔\quad 间隔\\\hline
间隔\texttt{\textbackslash qquad}间隔 & 间隔\qquad 间隔\\\hline
间隔$\sim$间隔 & 间隔~间隔\\\hline
间隔\texttt{\textbackslash ,}间隔 & 间隔\, 间隔\\\hline\hline
\texttt{\textbackslash hfill} & 水平填充\\\hline
\texttt{\textbackslash vfill} & 垂直填充\\\hline
\texttt{\textbackslash hspace\{5cm\}} & 水平空白5cm的距离\\\hline
\texttt{\textbackslash vspace\{5cm\}} & 垂直空白5cm的距离\\\hline
\end{tabular}
\end{center}
\end{table}
\LaTeX{}支持的单位有：\texttt{in}---英寸(inch)，\texttt{mm}---毫米，\texttt{cm}---厘米，\texttt{pt}---Points(大约1\slash 72~inch)，以及相对距离：\texttt{em}---当前字体的字符“\textbf{M}”的宽度，\texttt{ex}---当前字体的字符“\textbf{x}”的宽度，长度可以为负值，如~-1.5em~。\par
需要换行的地方请用“\textbackslash\textbackslash ”来控制，每段段末请加“\texttt{\textbackslash par}”来表示一个段落的结束。换行的段落结束的区别在于换行后不会有缩进，而段落结束的话，下一段会自动缩进两个汉字距离。
\section{章节标题以及条目的设置}
新的一章将在新的页面开始，使用\texttt{\textbackslash chapter}\{\LaTeX{}基础语法\}命令时，即可以得到~\hyperref[chapter-yufa]{第5章~\LaTeX{}基础语法}~这样的章名称。\par
同样的，还有命令\textbackslash section\{•\}、\textbackslash subsection\{•\}来定义节、条的标题。对于章节前面的“第1章”、“1.1”等标号，则不需要用户来编辑，\LaTeX{}会根据规则来进行编排。
\section{图形和多图排列}
在编写论文过程中，需要引用图形甚至多个图形时，可以使用\texttt{figure}环境，以代码\ref{code-fig}来说明如何使用此环境。\par
需要注意的是，figure为浮动环境，即\LaTeX{}在排版浮动环境时，会根据一定的美观规则来排版，并不是将将代码放在某处，图片就会出现在某处，如代码\ref{code-fig}中的\texttt{[htbp]}为浮动环境选项，\texttt{h}表示为here，\texttt{t}为top，\texttt{b}为bottom，\texttt{p}为page，优先顺序依次递减，\LaTeX{}先判断here（当前位置）是否可以放置此环境（代码\ref{code-fig}中为图片），否则依次判断top、bottom、page来放置此环境中的内容。
\begin{lstlisting}[label={code-fig},caption={使用figure环境},language={LaTeX}]
 如图 \ref{fig-buaamark}所示为北航校徽。
\begin{figure}[htbp]
\begin{center}
\includegraphics[width=90bp]{figure/buaamark.eps}
\caption{校徽\label{fig-buaamark}}
\end{center}
\end{figure}
\end{lstlisting}
代码中\textbackslash includegraphics后的[width=90bp]为对图片大小进行控制的代码，可以替换成为height,scale等。\textbackslash caption为图的标题，\textbackslash label为图的引用标号，在“如图”后引用即可得到“\ref{figbuaamark}”的效果。具体生成的效果如第\pageref{figbuaamark}页所示。\par
如图\ref{figbuaamark}所示为北航校徽。\par
\begin{figure}[htb]
\begin{center}
\includegraphics[width=90bp]{figure/buaamark.eps}
\caption{校徽\label{figbuaamark}}
\end{center}
\end{figure}
而对于需要两个或三个图形并排或者矩形排列的时候，则示例代码可以如\ref{code-subfigs}所示。

\begin{lstlisting}[caption={图形的排列},label={code-subfigs},language={LaTeX}]
\begin{figure}[htb]
\begin{center}
    \subfigure[ 并排小图 a]{\label{fig:minia}\includegraphics[width=.2\textwidth]{buaamark.eps}}\hspace{7em}
    \subfigure[ 并排小图 b]{\label{fig:minib}\includegraphics[width=.2\textwidth]{buaamark.eps}}
    \caption{ 多个图形并排的示例 }
    \label{subfig1}
\end{center}
\end{figure}

\begin{figure}[htb]
\begin{minipage}{0.5\textwidth}
    \centering
    \includegraphics[width=0.35\textwidth]{buaamark.eps}
    \caption{ 小图A}
    \label{minifiga}
\end{minipage}
\begin{minipage}{0.5\textwidth}
    \centering
    \includegraphics[width=0.35\textwidth]{buaamark.eps}
    \caption{ 小图B}
    \label{minifigb}
\end{minipage}
\end{figure}

\begin{figure}[htb]
\centering
    \begin{tabular}{cc}
    \subfigure[ 子图 A]{\label{suba}\includegraphics[width=.18\textwidth]{figure/buaamark.eps}}\hspace{7em} & \subfigure[ 子图 B]{\label{subB}\includegraphics[width=.18\textwidth]{figure/buaamark.eps}}\\
    \subfigure[ 子图 C]{\label{subc}\includegraphics[width=.18\textwidth]{figure/buaamark.eps}}\hspace{7em} & \subfigure[ 子图 D]{\label{subd}\includegraphics[width=.18\textwidth]{figure/buaamark.eps}}\\
    \end{tabular}
\caption{ 矩形的 subfig 排列}
\label{maxfig}
\end{figure}
\end{lstlisting}
而具体的效果如下所示。\par
图\ref{subfig1}为子图排列的形式，两个子图有各自的图题，并有一个共同的图题。\par
\begin{figure}[htb]
\begin{center}
    \subfigure[并排小图a]{\label{fig:minia}\includegraphics[width=.2\textwidth]{buaamark.eps}}\hspace{7em}
    \subfigure[并排小图b]{\label{fig:minib}\includegraphics[width=.2\textwidth]{buaamark.eps}}
    \caption{多个图形并排的示例}
    \label{subfig1}
\end{center}
\end{figure}
图\ref{minifiga}和图\ref{minifigb}为两个并排的图形，各自均有图题。\par
\begin{figure}[htb]
\begin{minipage}{0.5\textwidth}
    \centering
    \includegraphics[width=0.35\textwidth]{buaamark.eps}
    \caption{小图A}
    \label{minifiga}
\end{minipage}
\begin{minipage}{0.5\textwidth}
    \centering
    \includegraphics[width=0.35\textwidth]{buaamark.eps}
    \caption{小图B}
    \label{minifigb}
\end{minipage}
\end{figure}
图\ref{maxfig}所示为2x2排列的图形，使用\texttt{tabular}的表格环境嵌套在\texttt{figure}环境中来达到此效果。\par
\begin{figure}[htb]
\centering
    \begin{tabular}{cc}
    \subfigure[子图A]{\label{suba}\includegraphics[width=.18\textwidth]{figure/buaamark.eps}}\hspace{7em} & \subfigure[子图B]{\label{subB}\includegraphics[width=.18\textwidth]{figure/buaamark.eps}}\\
    \subfigure[子图C]{\label{subc}\includegraphics[width=.18\textwidth]{figure/buaamark.eps}}\hspace{7em} & \subfigure[子图D]{\label{subd}\includegraphics[width=.18\textwidth]{figure/buaamark.eps}}\\
    \end{tabular}
\caption{矩形的subfig排列}
\label{maxfig}
\end{figure}
到此，图形的排列基本就OK了。一般只会用到图\ref{figbuaamark}这样的图形环境就可以了，后面的太复杂，一般情况下不推荐使用。如果对图形的编排的兴趣，可以google一下《\LaTeX{}插图指南》一书。
\section{表格和长表格}
    如果您要排版的表格长度超过一页，那么推荐使用 \textsf{longtable} 或者 \textsf{supertabular} 
宏包，模板对 \textsf{longtable} 进行了相应的设置，所以用起来可能简单一些。
表~\ref{example:longtable} 就是 \textsf{longtable} 的简单示例。
\begin{center}
\begin{longtable}{|l|l|l|}
\caption{长表格示例} \label{example:longtable} \\

\hline \multicolumn{1}{|c|}{\textbf{Time (s)}} & \multicolumn{1}{c|}{\textbf{Triple chosen}} & \multicolumn{1}{c|}{\textbf{Other feasible triples}} \\ \hline 
\endfirsthead

\multicolumn{3}{l}%
{{\bfseries \tablename\ \thetable{} --接\,上\,页}} \\
\hline \multicolumn{1}{|c|}{\textbf{Time (s)}} &
\multicolumn{1}{c|}{\textbf{Triple chosen}} &
\multicolumn{1}{c|}{\textbf{Other feasible triples}} \\ \hline 
\endhead

\hline \multicolumn{3}{|r|}{{接\,下\,页}} \\ \hline
\endfoot

\hline \hline
\endlastfoot

0 & (1, 11, 13725) & (1, 12, 10980), (1, 13, 8235), (2, 2, 0), (3, 1, 0) \\
2745 & (1, 12, 10980) & (1, 13, 8235), (2, 2, 0), (2, 3, 0), (3, 1, 0) \\
5490 & (1, 12, 13725) & (2, 2, 2745), (2, 3, 0), (3, 1, 0) \\
8235 & (1, 12, 16470) & (1, 13, 13725), (2, 2, 2745), (2, 3, 0), (3, 1, 0) \\
10980 & (1, 12, 16470) & (1, 13, 13725), (2, 2, 2745), (2, 3, 0), (3, 1, 0) \\
13725 & (1, 12, 16470) & (1, 13, 13725), (2, 2, 2745), (2, 3, 0), (3, 1, 0) \\
16470 & (1, 13, 16470) & (2, 2, 2745), (2, 3, 0), (3, 1, 0) \\
19215 & (1, 12, 16470) & (1, 13, 13725), (2, 2, 2745), (2, 3, 0), (3, 1, 0) \\
21960 & (1, 12, 16470) & (1, 13, 13725), (2, 2, 2745), (2, 3, 0), (3, 1, 0) \\
24705 & (1, 12, 16470) & (1, 13, 13725), (2, 2, 2745), (2, 3, 0), (3, 1, 0) \\
27450 & (1, 12, 16470) & (1, 13, 13725), (2, 2, 2745), (2, 3, 0), (3, 1, 0) \\
30195 & (2, 2, 2745) & (2, 3, 0), (3, 1, 0) \\
32940 & (1, 13, 16470) & (2, 2, 2745), (2, 3, 0), (3, 1, 0) \\
35685 & (1, 13, 13725) & (2, 2, 2745), (2, 3, 0), (3, 1, 0) \\
38430 & (1, 13, 10980) & (2, 2, 2745), (2, 3, 0), (3, 1, 0) \\
41175 & (1, 12, 13725) & (1, 13, 10980), (2, 2, 2745), (2, 3, 0), (3, 1, 0) \\
43920 & (1, 13, 10980) & (2, 2, 2745), (2, 3, 0), (3, 1, 0) \\
150975 & (1, 13, 16470) & (2, 2, 2745), (2, 3, 0), (3, 1, 0) \\
153720 & (1, 12, 13725) & (2, 2, 2745), (2, 3, 0), (3, 1, 0) \\
156465 & (1, 13, 13725) & (2, 2, 2745), (2, 3, 0), (3, 1, 0) \\
159210 & (1, 13, 13725) & (2, 2, 2745), (2, 3, 0), (3, 1, 0) \\
161955 & (1, 13, 16470) & (2, 2, 2745), (2, 3, 0), (3, 1, 0) \\
164700 & (1, 13, 13725) & (2, 2, 2745), (2, 3, 0), (3, 1, 0) \\
\end{longtable}
\end{center}
\section{公式和定理定义等}
    \section{使用~listing~显示代码}
原先~ctexbook~文档类和~listings~宏包配合使用时，代码在换页时会出现莫名其妙的错误，后来经高人指点，顺利解决了。
感兴趣的话，可以看看\href{http://bbs.ctex.org/viewthread.php?tid=53451}{这里}。
这里给使用~listings~宏包插入源代码的例子，这里是一段C代码。
另外，listings宏包真可谓博大精深，可以实现各种复杂、漂亮的效果，想要进一步学习的同学，可以参考
\href{http://mirror.ctan.org/macros/latex/contrib/listings/listings.pdf}{listings宏包手册\footnote{此节完全参考武汉理工大学论文模板}}。